% Transcription — fonds Alexandre Grothendieck, shelfmark 1, batch 6 (pages 101-120).
% Established from the facsimile archives/batches/1/batch-06.pdf (a mirror of
% https://grothendieck.umontpellier.fr/1.pdf), per the skill
% transcribe-grothendieck.
%
% Pass: Fable 5.1 (claude-fable-5-1), 2026-09-10 — first pass, unchecked against
% the pages by a human.
%
% Nineteen of the twenty pages are transcribed; page 105 is blank. \page{}
% numbers the position in the folder, as batches 1 to 5 did: the archivists'
% pencil numbers run two ahead, so page 101 here is pencilled 103 and page
% 120 is pencilled 122.
%
% The batch falls in two parts, in two hands.
%
%   Pages 101-104 — « Espaces de Schatten », in the fast blue-black hand of
%     batches 3 to 5: a plan in two points (① Normes de Schatten, ② Espaces
%     de Schatten) on pages 101-102, then a recapitulating « Théorème ① » in
%     five points on pages 103-104, whose formulas (18) to (27) continue the
%     numbering (1) to (17) of the n° 5 transcribed at pages 93-100 of batch
%     5: (18) N(u) = N((ρ_i(u))), (19) N(u) = Sup |Tr vu| over N°(v) ≤ 1, the
%     spaces L^N(E,F) and L_N(E,F), (20) N(u) = N(u*), (21) N(uv) ≤ ‖u‖N(v),
%     (22)-(23) the multiplicativity N(vu) ≤ N'(u)N''(v), (24)-(25) ‖vu‖_1
%     and |Tr vu| ≤ N(u)N°(v), (26) the pairing ⟨u,v⟩ = Tr vu identifying
%     (L_N(E,F))' with L^{N°}(F,E) and the reflexivity equivalences, (27)
%     N((λ_i(u))) ≤ N(u), and the proof of points 4 and 5.
%   Pages 106-120 — older, smaller sheets in a slower, more legible hand,
%     with printed versos showing through; nothing on them is dated.
%     Pages 106-111: « Décomposition en facteurs des fonctions entières » —
%     Th. 1 (Weierstrass) with (1)-(2), genre with (3)-(4), order, Th. 2
%     (Hadamard-Borel) with (5), the frame ω−1 ≤ g ≤ ω, Th. 3 (Hadamard) on
%     the coefficients, and a Proposition: an entire function whose
%     coefficients are dominated by those of a genre-zero function and whose
%     zeros satisfy Σ1/|z_i| < ∞ is of genre zero, proved on a sector with
%     G(z) = g(z)g(ωz)g(ω²z), ω = e^{2iπ/3}; the proof ends in the margin.
%     Pages 112-114: the exponent of convergence ρ of a decreasing sequence,
%     the lemma α_n = o(n^{-1/s}), ρ as the order of ∏(1+α_i z) through the
%     case α_n = n^{-β}, f and F = ∏(1+|λ_i|z) having the same order, and a
%     half page of very faint pencil on orders of sequences in a Banach space.
%     Pages 115-118: the counting function n(r), (1) and (1') the Stieltjes
%     summation Σ r_ν^{-α} = n(b−0)/b^α − n(a+0)/a^α + α∫n(r)r^{-α-1}dr, the
%     proposition Σ r_ν^{-α} < ∞ ⟺ ∫n(r)/r^{α+1} < ∞ with n(r)/r^α → 0, its
%     form for Σλ(r_ν), and the Jensen-type identity ∫n(r)μ(r)dr =
%     [rμV]_a^b − ∫V d(rμ) with V(r) = (1/2π)∫log|f(re^{iφ})|dφ.
%     Pages 119-120: « Sur une certaine fonction holomorphe sur l'espace de
%     Banach ℓ¹ » — α_n(λ) = Σ λ_{i_1}⋯λ_{i_n} ≤ ‖λ‖_1^n/n!, the Proposition
%     that Σ n!α_n(λ) converges absolutely with S(λ) holomorphic near 0, the
%     Corollaire ⁿ√|α_n(λ)| = o(1/n), and the start of the proof through
%     Cauchy's inequality with R = (n+1)/ρ_{n+1}; it breaks off at the foot
%     of page 120 and continues in batch 7.
%
% Cancellations: pages 101, 102, 103, 104, 106, 107, 108, 109, 110, 111,
% 112, 113, 114, 115, 116, 117, 118 and 120 carry struck words or lines,
% transcribed as \struck{} with one \note{} per block where the block cannot
% be read. Page 101 draws a large freehand frame around its two definitions
% of N(u). Red pencil underlines on page 102 and a red question mark in the
% margin of page 103. The sideways note in the left margin of page 111 is
% transcribed as \marginal{}; page 117 carries two oblique words and a
% freehand circle in its margin. The lower half of page 114 is pencil, very
% faint, and is not recovered.
%
% Notation: as in batch 5 — ρ_i(u) the singular values, λ_i(u) the
% eigenvalues, N a Schatten norm and N° its polar, ℓ^N, ℓ_N, L^N(E,F),
% L_N(E,F), c_0, R^(N) set \mathbf{R}^{(\mathbf{N})} (page 101 also writes
% C^N and C^(N), kept), E'⊗F the finite-rank operators. In the second part:
% M(r) = Sup_{|z|=r}|f(z)|, the order written ρ on page 107 and ω from page
% 108 on, both kept; « C^{te} » is his abbreviation for a constant, set
% C^{\mathrm{te}}; n(r) the counting function; Sup, Max, lim sup set as he
% writes them, the last as \limsup.
%
% The dating is the inventory's own for the shelfmark. Its group,
% « MATHÉMATIQUES / RECHERCHE / Travaux », carries no date.
\documentclass[11pt,a4paper]{article}
% Le préambule commun à toutes les transcriptions du fonds Grothendieck.
%
% Il tient en une idée : **l'incertitude se compose, elle ne se résout pas.**
% Une transcription qui lisse un mot douteux a détruit la seule chose pour
% laquelle on la faisait — un lecteur doit pouvoir savoir, à chaque endroit,
% ce qui était sur le feuillet et ce qui a été ajouté. Les six macros qui
% suivent sont donc l'appareil critique, et non de la décoration.
%
% Les mêmes commandes sont reconnues par `scripts/render.mjs`, qui produit la
% vue de lecture du site. Une macro ajoutée ici doit l'être aussi là-bas,
% faute de quoi le rendu échouera bruyamment — ce qui est voulu : un rendu
% silencieusement faux est pire qu'un rendu absent.

\ProvidesPackage{grothendieck}[2026/08/08 Transcriptions du fonds Grothendieck]

\usepackage{amsmath,amssymb,amsthm}
\usepackage{mathrsfs}
\usepackage{stmaryrd}
% Les diagrammes commutatifs sont partout dans ce fonds : c'est la langue
% même de la géométrie algébrique de Grothendieck, et une transcription qui
% les rendrait en prose ne transcrirait rien.
\usepackage{tikz-cd}
% Français par défaut — c'est la langue des feuillets — mais l'anglais est
% chargé aussi : la traduction et le résumé sont des documents anglais, et la
% typographie française (espaces avant les deux-points) y serait une faute.
% Ces documents basculent par \selectlanguage{english} en tête de corps.
\usepackage[english,french]{babel}
\usepackage{fontspec}
\usepackage{geometry}
\geometry{a4paper,margin=25mm,marginparwidth=28mm}
\usepackage{marginnote}
\usepackage{xcolor}
% ulem, not soul. `soul`'s \st and \ul are built on a character-by-character
% scan that XeLaTeX's font machinery breaks: the document fails with an
% undefined control sequence rather than degrading. ulem does the same two jobs
% and is Unicode-engine safe. `normalem` stops it hijacking \emph, which the
% transcriptions use for Grothendieck's own emphasis.
\usepackage[normalem]{ulem}
\usepackage{hyperref}
\hypersetup{colorlinks=true,linkcolor=black,urlcolor=black,pdfborder={0 0 0}}

% --- Métadonnées du lot -----------------------------------------------------
% Reprises telles quelles par le script de rendu, qui les affiche en tête de
% la vue de lecture. Elles fixent ce que le fichier prétend couvrir : sans
% elles, une transcription flotte sans point d'attache dans un fonds de seize
% mille feuillets.

\newcommand{\folder}[1]{\gdef\grothFolder{#1}}
\newcommand{\batch}[1]{\gdef\grothBatch{#1}}
\newcommand{\leaves}[2]{\gdef\grothFirst{#1}\gdef\grothLast{#2}}
\newcommand{\foldertitle}[1]{\gdef\grothFoldertitle{#1}}
\gdef\grothFolder{?}\gdef\grothBatch{?}\gdef\grothFirst{?}\gdef\grothLast{?}\gdef\grothFoldertitle{}

% --- Le résumé --------------------------------------------------------------
%
% Il ouvre la lecture modernisée, sous le titre « Résumé », et il est composé
% en petit. Ce n'est pas un ornement : c'est le seul passage du document qui
% ne soit pas des mathématiques, et un lecteur doit pouvoir le reconnaître
% avant de l'avoir lu — puis le sauter s'il connaît déjà le sujet, ou s'y
% arrêter s'il ne le connaît pas. Le filet qui le suit marque la reprise du
% corps.

\newenvironment{resume}
  {\small\setlength{\parskip}{0.45em}}
  {\par\medskip\hrule\medskip}

% --- Mots-clés ---------------------------------------------------------------
%
% En anglais, en clôture du résumé de la lecture modernisée : le vocabulaire
% moderne sous lequel chercher ce que les feuillets construisent. Le manifeste
% du site (`npm run manifest`) les extrait pour étiqueter la cote entière —
% cette ligne est la source unique des tags, il n'y a pas de fichier à côté.
\newcommand{\keywords}[1]{\par\smallskip\noindent
  {\footnotesize\textsc{Keywords} --- \textit{#1}}\par}

% --- L'appareil critique ----------------------------------------------------

% Le numéro de feuillet, en marge. C'est l'ancre qui permet de régler un
% désaccord sur une formule en regardant *un* feuillet plutôt qu'une chemise
% de six cents. Le site s'en sert aussi pour tourner les pages du fac-similé
% quand on fait défiler la transcription.
\newcommand{\leaf}[1]{%
  \marginnote{\footnotesize\color{gray!70}\textsf{#1}}%
  \label{leaf:#1}\ignorespaces}

% Illisible : on ne devine pas. Un mot inventé qui se lit comme les autres est
% le pire résultat possible.
\newcommand{\ill}{\textcolor{red!60!black}{[\,\ldots\,]}}

% Lecture douteuse : le mot est donné, et signalé comme incertain.
\newcommand{\uncertain}[1]{\uline{#1}}

% Ajout de l'éditeur — une lettre manquante, un mot sous-entendu.
\newcommand{\add}[1]{\textcolor{blue!50!black}{[#1]}}

% Biffé par Grothendieck lui-même. Ce qu'il a rayé fait partie du document :
% une rature dit souvent où la pensée a changé de direction.
\newcommand{\struck}[1]{\sout{#1}}

% Note du transcripteur, sur l'état matériel du feuillet.
\newcommand{\note}[1]{\par\smallskip{\small\color{gray!60!black}\textit{[#1]}}\par\smallskip}

% Note marginale de Grothendieck, distincte du corps du texte.
\newcommand{\marginal}[1]{\par\smallskip\noindent\hspace{1em}%
  {\small\color{gray!60!black}#1}\par\smallskip}

% --- Titre ------------------------------------------------------------------

\AtBeginDocument{%
  \begin{center}
    {\large\bfseries Fonds Alexandre Grothendieck --- cote n\textsuperscript{o}~\grothFolder}\\[2pt]
    {\itshape\grothFoldertitle}\\[4pt]
    {\small feuillets \grothFirst--\grothLast\ (lot \grothBatch)}\\[2pt]
    {\footnotesize\color{gray!60!black}Transcription établie d'après le fac-similé numérisé
     par l'Université de Montpellier.\\
     Les lectures incertaines sont soulignées, les illisibles notées [\,\ldots\,],
     les ajouts entre crochets.}
  \end{center}
  \vspace{1em}}

\endinput


\folder{1}
\batch{6}
\pages{101}{120}
\dating{1953}
\watermark{Édition de démonstration}
\foldertitle{Analyse fonctionnelle : notes manuscrites (s.d.), lettre (1953)}

\begin{document}

\section{Espaces de Schatten}

\note{Le titre est de lui, encadré, en tête de la page 101. Les pages 101 à 104 sont un plan en deux points, puis un énoncé récapitulatif, « Théorème 1 », en cinq points ; les formules y sont numérotées de (18) à (27), à la suite des (1) à (17) du n° 5 transcrit aux pages 93 à 100 du lot précédent. Les numéros des points sont cerclés sur la page et rendus ici par 1., 2., …}

\page{101}
\textbf{1. Normes de Schatten.} Définition. \emph{Prop} Une norme de Sch\add{atten} $N(x)$ est \uncertain{fonction} \uncertain{croissante} de $|x|$. Exemples : les $N_p$. Norme de Schatten \uncertain{définie} par une norme de \uncertain{Schatten} sur $\mathbf{R}^n$. Définition $N(x)$ pour $x \in \mathbf{C}^{\mathbf{N}}$. $\ell^N \ni x \iff N(x)$ fini, $N$ est une norme sur $\ell^N$. Espace $\ell_N$, adhérence de $\mathbf{C}^{(\mathbf{N})}$ dans $\ell^N$. Inclusions pour $N' \leq N''$, \uncertain{en} particulier : $\ell^1 \subset \ell^N \subset \ell^\infty$, $\ell^1 \subset \ell_N \subset c_0$. \emph{Prop} $\ell_N \subset c_0$ sauf si $N$ est équivalente à $N_\infty$. Polarité, $(N^{\circ})^{\circ} = N^{\circ}$ \note{ainsi sur la page ; on attend $(N^{\circ})^{\circ} = N$} ; \emph{Prop} le dual de $\ell_N$ est $\ell^{N^{\circ}}$. \emph{Corollaire} $\ell^N$ est \uncertain{complet}, $\ell_N$ est complet. \emph{Prop} $\ell_N$ réflexif $\iff \ell^N$ réflexif $\iff \ell_{N^{\circ}}$ réflexif $\iff \ell^{N^{\circ}}$ réflexif $\iff \ell_N = \ell^N$ et $\ell_{N^{\circ}} = \ell^{N^{\circ}}$. \struck{Exemples} Cas des $N_p$ : \ill{} $(N_p)^{\circ} = N_{p'}$. On retrouve la réflexivité des $\ell^p$ ($1 < p < +\infty$).

\textbf{2. Espaces de Schatten.} Définition si $u \in L_0(E,F)$, \struck{$\Longrightarrow$} $N(u) = N((\rho_i(u)))$ ; \ill{} $u \in L^N(E,F) \iff N(u) < +\infty$, \uncertain{norme} sur $L^N(E,F)$, $L_N(E,F)$ adhérence de $E' \otimes F$ ($N$ non équivalente à $N_\infty$). $N(u) = N(u^{*})$, $N(vu) \leq \|v\|\,N(u)$. Définition (si $N$ \uncertain{équivalente} à $N_\infty$) $N(u) = \operatorname{Sup}_{N^{\circ}(v) \leq 1,\ v \in F' \otimes E} |\operatorname{Tr} vu|$, si $u \in L_0(E,F)$, on retrouve $N(u) = N((\rho_i(u)))$, i.e. \ill{}
\[
  N((\rho_i(u))) = \operatorname{Sup}_{v \in F' \otimes E,\ N^{\circ}(v) \leq 1} |\operatorname{Tr} vu| .
\]
\note{les deux définitions et cette formule sont réunies par un grand cadre à main levée, avec des traits de renvoi vers le bas de la page}

Définition, si $u \in L_0(E,F)$, posons $N(u) = N((\rho_i(u)))$ ; $N(u) = \operatorname{Sup}_{v \in F' \otimes E,\ N^{\circ}(v) \leq 1} |\operatorname{Tr} vu|$ \note{la condition $N^{\circ}(v) \leq 1$ sous le Sup est surchargée, « $N(v)$ » récrit au-dessous}. Définition, pour $u \in L(E,F)$, $N(u)$ par la formule ci-dessus. Soit \struck{$N(u) = \operatorname{Sup} \operatorname{Tr} vu$} $L^N(E,F) \ni u \iff N(u) < +\infty$, $N(u)$ une norme sur $L^N(E,F)$. \struck{et on a} Soit $L_N(E,F)$ l'adhérence de $E' \otimes F$ dans $L^N(E,F)$. \emph{Prop} $u \in L_N(E,F) \iff u \in L_0(E,F)$ et $(\rho_i(u)) \in \ell_N$. Inclusions pour $N' \leq N''$ $\Longrightarrow$ $L^{N''}(E,F) \subset L^{N'}(E,F)$ et $L_{N''}(E,F) \subset L_{N'}(E,F)$, en particulier $L^1(E,F) \subset L^N(E,F) \subset L(E,F)$, $L^1(E,F) \subset L_N(E,F) \subset L_0(E,F)$. \struck{\emph{Prop} Soit $N$ \ill{} $N_\infty$, alors} \emph{Prop} $N(u) = N(u^{*})$, $N(vu) \leq \|v\|\,N(u)$, $N(uv) \leq \|v\|\,N(u)$ \note{les trois formules sont encadrées}. \struck{\emph{Corollaire} \ill{}} (Pour \uncertain{prouver} \ill{}, il suffit de le vérifier pour $u \in L_0(E,F)$, où \add{cela} résulte de $\rho_i(vu) \leq \|v\|\,\rho_i(u)$ \ill{}.)

\page{102}
\marginal{\uncertain{Aussi} le \uncertain{dual} \ill{} $L^N$ \ill{}} \note{trois lignes courtes dans le coin supérieur gauche}

\emph{Corollaire} $N(u) = N(|u|)$ ; si $u = h + ik$, $N(h), N(k) \leq N(u) \leq N(h) + N(k)$, donc $u \in L^N(E,F) \iff h, k \in L^N(E,F)$.

\emph{Prop} $L^N(E,F) \subset L_0(E,F)$ si $N$ non équivalente à $N_\infty$ (sinon $L^N(E,F) = L(E,F)$). En effet, si $u \in L^N(E,F)$, \struck{\ill{}} alors $|u| \in L^N(E)$, \uncertain{prouvons} $|u|$ compact (d'où $u$ compact). Sinon, il existerait un op\add{érateur} hermitien $v \in L(E)$ \uncertain{tel} \uncertain{que} $v|u|$ soit un projecteur hermitien \uncertain{infini}, or on aurait $v|u| \in L^N(E)$. C'est impossible, car on \struck{aurait $(\rho_i) \in \ell^N$} aurait $\sum |x_i| \leq M$ pour $N^{\circ}((x_i)) \leq 1$, donc $N^{\circ}$ équivalent à $N_1$, ce qui est absurde. \note{de « C'est impossible » à « absurde », le passage est encadré}

\emph{Corollaire} $L^{N_p}(E,F) = L^p(E,F)$ pour $1 \leq p \leq +\infty$, \note{en interligne : « si $p < +\infty$, on a »} $L_{N_p}(E,F) = L^p(E,F)$ (car $\ell_N = \ell^N \Longrightarrow L_N(E,F) = L^N(E,F)$).

\emph{Prop} $u \in L^N(E,F)$, $v \in L^{N^{\circ}}(F,E) \Longrightarrow vu \in L^1(E,G)$ \note{ainsi sur la page ; on attend $L^1(E,E)$}. D'où accouplement $\operatorname{Tr} vu$ entre $L^N(E,F)$ et $L^{N^{\circ}}(F,E)$. \struck{\ill{}} \note{deux lignes biffées} \emph{Prop} le dual de $L_N(E,F)$ est $L^{N}(E,F)$ \note{ainsi sur la page : ni le « ° » ni l'ordre $(F,E)$ attendus ne sont lisibles ; le point 4 du théorème 1, page 104, donne $L^{N^{\circ}}(F,E)$}. \emph{Corollaire 1} $L^N(E,F)$ et $L_N(E,F)$ sont complets. \emph{Corollaire 2} $L_N(E,F)$ réflexif $\Longleftarrow$ $N$ réflexif \note{souligné au crayon rouge}. Réciproque vraie si $E$ et $F$ de dimension infinie. En tous cas, $L_N(E,F)$ réfl\add{exif} $\iff L^N(E,F)$ réflexif $\iff$ \ill{} \note{souligné au crayon rouge}. \emph{Corollaire 3} Si $1 < p < +\infty$, alors $L^p(E,F)$ est un espace réflexif ayant pour dual $L^{p'}(F,E)$.

\emph{Prop} Si $N, N', N''$ sont des normes de Schatten telles que $N(xy) \leq N'(x)\,N''(y)$, alors $u \in L^{N'}(E,F)$, $v \in L^{N''}(F,G)$ implique $vu \in L^N(E,F)$ \note{ainsi sur la page ; $(E,G)$ attendu}, et $N(vu) \leq N'(u)\,N''(v)$. \note{en interligne, cerclé : « \emph{Corollaire} si on a \ill{} $N$, $N''$ \ill{} »} (Pour \uncertain{prouver} la proposition, \ill{} par le deuxième th\add{éorème} de convexité si $v$ et $u$ sont \ill{} ; \struck{\ill{}} \note{deux lignes biffées} si $N(xyz) \leq N'(x)\,N''(y)\,N'''(z)$, alors même formule vraie \ill{} \struck{$N(wvu) \leq N'(u)\,N''(v)\,N'''(w)$}. $N(u) = \lim N(pup) = \lim_p N(up) = \lim_p N(pu)$ ; $p$ projecteur orth\add{ogonal} de rang fini, croissant \uncertain{vers} l'identité. (Alors $N(pvup) = N((pv)(up)) \leq N''(pv)\,N'(up) \leq N''(v)\,N'(u)$.) \ill{} en dimension finie \ill{}.)

\emph{Prop} $N((\lambda_i(u))) \leq N(u)$.

\emph{Prop} (caractérisation \uncertain{axiomatique} des espaces de Schatten\ill{}).

\page{103}
\emph{Théorème 1} \note{le « 1 » est cerclé} Soit $N$ une norme de Schatten sur $\mathbf{R}^{(\mathbf{N})}$, $E$ et $F$ deux espaces de \struck{Banach} \note{en interligne : « Hilbert »}, et pour $u \in L_0(E,F)$, posons
\[
  \text{(18)}\qquad N(u) = N((\rho_i(u)))
\]
et pour $u \in L(E,F)$
\[
  \text{(19)}\qquad N(u) = \operatorname{Sup}_{v \in F' \otimes E,\ N^{\circ}(v) \leq 1} |\operatorname{Tr} vu|
\]
\uncertain{Ces} \uncertain{deux} \uncertain{expressions} coïncident \ill{} sur $L_0(E,F)$. Soit $L^N(E,F)$ \struck{l'ensemble} le sous-espace de $L(E,F)$ formé des $u$ tels que $N(u) < +\infty$, alors $N$ est une norme sur $L^N(E,F)$, qui en fait un espace de Banach. \struck{L'adhérence de $E' \otimes F$ dans $L^N(E,F)$ \ill{}} \struck{\ill{}} \note{deux passages biffés, de deux lignes chacun} \struck{Pour que $u \in L_0(E,F)$ \ill{}} \uncertain{Si} $u \in L_0(E,F)$, $u \in L^N(E,F)$ si et seulement si $(\rho_i(u)) \in \ell^N$, et alors on a (18). L'adhérence $L_N(E,F)$ de $E' \otimes F$ dans $L^N(E,F)$ est l'ensemble des $u \in L_0(E,F)$ tels que $(\rho_i(u)) \in \ell_N$.

2. On a
\[
  \text{(20)}\qquad N(u) = N(u^{*})
\]
\note{formule encadrée} et par suite $u \in L^N(E,F) \iff u^{*} \in L^N(F,E)$, $u \in L_N(E,F) \iff u^{*} \in L_N(F,E)$. On a
\[
  \text{(21)}\qquad N(uv) \leq \|u\|\,N(v), \qquad N(vu) \leq \|v\|\,N(u)
\]
si $u \in L(E,F)$, $v \in L(F,G)$ \note{ainsi sur la page, les rôles de $u$ et $v$ ne s'accordant pas d'une inégalité à l'autre ; en marge, biffé : « sur $L^N(E,F)$ », avec un point d'interrogation au crayon rouge}. \uncertain{Donc}, \ill{} : \struck{\ill{}} \uncertain{de} type $\ell^N$ (resp. $\ell_N$) \ill{} \uncertain{injection}, on obtient un espace du type $\ell^N$ (resp. $\ell_N$). \note{en interligne : « $\Longrightarrow$ $N(u)$ est fonction croissante de $u$ sur l'ens\add{emble} des \uncertain{hermitiens} positifs dans $L_0(E)$ »}

3. \struck{4.} Si $N, N', N''$ sont des normes de Schatten telles que
\[
  \text{(22)}\qquad N(xy) \leq N'(x)\,N''(y)
\]
\note{formule encadrée} pour $x, y \in \mathbf{R}^{(\mathbf{N})}$, alors on a, pour $u \in L(E,F)$, $v \in L(F,G)$,
\[
  \text{(23)}\qquad N(vu) \leq N'(u)\,N''(v)
\]
donc si $u \in L^{N'}(E,F)$, $v \in L^{N''}(E,F)$, alors $vu \in L^N(E,F)$ \note{ainsi sur la page ; $(F,G)$ et $(E,G)$ attendus}, et si de plus $u \in L_{N'}(E,F)$ ou $v \in L_{N''}(E,F)$, on aura $vu \in L_N(E,F)$. --- En particulier, on aura, pour $u \in L(E,F)$, $v \in L(F,E)$
\[
  \text{(24)}\qquad \|vu\|_1 \leq N(u)\,N^{\circ}(v)
\]
donc si $u \in L^N(E,F)$, $v \in L^{N^{\circ}}(F,E)$, on a $vu \in L^1(E,E)$, et
\[
  \text{(25)}\qquad |\operatorname{Tr} vu| \leq N(u)\,N^{\circ}(v) .
\]

\page{104}
4. Par l'accouplement
\[
  \text{(26)}\qquad \langle u, v \rangle = \operatorname{Tr} vu
\]
le dual de $L_N(E,F)$ s'identifie à $L^{N^{\circ}}(F,E)$. Si $E$ et $F$ sont de dimension infinie, \struck{les} \note{en interligne : « alors »} \uncertain{les} conditions suivantes sont équivalentes : a) $L_N(E,F)$ est réflexif \struck{$L_N(E,F) = L^N(E,F)$ et $L_{N^{\circ}}(F,E) = L^{N^{\circ}}(F,E)$} b) $L^N(E,F)$ est réflexif c) $\ell_N$ est réflexif \struck{\ill{}} d) $\ell^N$ est réflexif ; e) $\ell_{N^{\circ}}$ est réflexif et f) $\ell^{N^{\circ}}$ est réflexif ; d) $L_N(E,F) = L^N(E,F)$ et $L_{N^{\circ}}(F,E) = L^{N^{\circ}}(F,E)$. \note{la liste est remaniée sur la page : les lettres sont surchargées, « d) » apparaît deux fois, et des traits relient les conditions entre elles ; l'ordre donné ici suit la ligne}

5. \struck{Remarque} Soit $u \in L_0(E,F)$ \struck{si $u \in L^N(E,F)$ alors $(\lambda_i(u)) \in \ell^N$}, \uncertain{alors}
\[
  \text{(27)}\qquad N((\lambda_i(u))) \leq N(u)
\]
\note{« $N((\rho_i(u)))$ » biffé entre les deux membres} Si $u \in L^N(E,F)$, alors $(\lambda_i(u)) \in \ell^N$ \note{« $\ell_N$ » biffé, « $\ell^N$ » récrit}.

Tout est démontré, sauf 4. et 5. Pour 4, on note : si $E$ et $F$ sont de dimension infinie, alors $\ell^N$ (resp. $\ell_N$) est isomorphe à un sous-espace normé de $L^N(E,F)$ (resp. $L_N(E,F)$). (On choisit $(e_i)$, $(f_i)$ systèmes orthonormaux infinis, et $(\lambda_i) \mapsto \sum \lambda_i\,\bar{e}_i \otimes f_i$ de $\ell^N$ dans $L^N(E,F)$.) Donc si $L^N(E,F)$ est réflexif, $L_N(E,F)$ est réflexif, et alors $\ell_N$ est réflexif, \struck{\ill{}} donc (voir \ill{}) $\ell_N = \ell^N$, $\ell_{N^{\circ}} = \ell^{N^{\circ}}$, donc $L^N(E,F) = L_N(E,F)$ (car \ill{} si $N \neq N_\infty$), donc $L^N(E,F)$ est réflexif, d'où l'équivalence de a), b), c), et le fait que ces conditions impliquent d) (donc aussi \uncertain{la} \uncertain{première} \uncertain{relation}, \ill{} $\ell_{N^{\circ}}$ est aussi réflexif). Enfin, si \ill{} le dual de $L_N(E,F)$ est $L_{N^{\circ}}(F,E)$, \struck{d'où $\ell_N = \ell^N$ et $\ell_{N^{\circ}} = \ell^{N^{\circ}}$ (voir \ill{})} \struck{\ill{} $\ell_N$ réflexif (voir \ill{})} d'où \ill{} $L_N(E,F)$, donc $L_N(E,F)$ est réflexif. 5 : (27) \struck{résulte immédiatement} est \uncertain{le} corollaire 5 du \uncertain{premier} th\add{éorème} de convexité.

\section{Décomposition en facteurs des fonctions entières}

\note{Le titre est de lui, souligné, en tête de la page 106. Les pages 106 à 120 sont d'autres feuillets, d'une écriture plus posée, dont les versos imprimés transparaissent ; la page 105 est blanche. Les formules sont numérotées de (1) à (5) aux pages 106 à 108.}

\page{106}
\emph{Théorème 1} (Weierstrass) Soit $f(z)$ fonction entière, $(a_n)$ la suite de ses zéros, rangés par ordre de modules croissants. Supposons $f(0) \neq 0$. Il est possible de trouver une suite croissante de nombres entiers $\geq 0$ $p_n$, et une fonction entière $g(z)$, tels que l'on ait
\[
  \text{(1)}\qquad f(z) = e^{g(z)} \prod_n \Bigl(1 - \frac{z}{a_n}\Bigr)\, e^{\frac{z}{a_n} + \frac{z^2}{2a_n^2} + \cdots + \frac{z^{p_n}}{p_n a_n^{p_n}}}
\]
le produit étant uniformément absolument convergent quand $z$ parcourt n'importe quelle région bornée. Si \struck{de plus} une suite d'entiers $p_n$ vérifie $\sum_n |r/a_n|^{p_n+1} < +\infty$ quel que soit $r$, alors on peut employer cette suite dans (1). En particulier, supposons $k$ un entier tel que $\sum 1/|a_n|^{k+1} < +\infty$, alors on peut prendre $p_n = p$ pour tout $n$, avec $p \leq k$ \note{ainsi lu ; (2) prend $p = k$} : d'où
\[
  \text{(2)}\qquad f(z) = e^{g(z)} \prod_n \Bigl(1 - \frac{z}{a_n}\Bigr)\, e^{\frac{z}{a_n} + \cdots + \frac{z^k}{k a_n^k}}
\]

\page{107}
\emph{Genre d'une fonction entière.} Il arrive que l'on ait $\sum 1/|a_n|^{k+1} < +\infty$ pour quelque $k$ (soit $k$ le plus petit entier qui y satisfasse) et de plus que dans la représentation (2), $g(z)$ soit un polynôme --- d'ailleurs déterminé à une constante près. Soit $r$ son degré ; le plus \struck{petit} grand des nombres $r$, $k$ s'appelle le genre de $f(z)$. C'est un nombre entier $\geq 0$. Dire que $f(z)$ est de genre $\leq r$, \uncertain{signifie} : on a la formule
\[
  \text{(3)}\qquad f(z) = e^{\alpha_0 + \alpha_1 z + \cdots + \alpha_r z^r} \prod_n \Bigl(1 - \frac{z}{a_n}\Bigr)\, e^{\frac{z}{a_n} + \cdots + \frac{z^r}{r a_n^r}}
\]
Dire que $f(z)$ est de genre zéro, signifie que $f(z)$ est de la forme
\[
  \text{(4)}\qquad f(z) = C^{\mathrm{te}} \prod_n \Bigl(1 - \frac{z}{a_n}\Bigr), \quad \text{avec } \sum_n \frac{1}{|a_n|} < +\infty .
\]

\emph{Ordre d'une fonction entière.} Posons $M(r) = \operatorname{Sup}_{|z| = r} |f(z)|$. Puis $\rho = \limsup \dfrac{\log\log M(r)}{\log r}$ $(0 \leq \rho \leq +\infty)$, $\rho$ est dit l'ordre de $f(z)$, et $f(z)$ est dite d'ordre fini si $\rho < +\infty$. Cela signifie qu'il existe $k > 0$ tel que \ill{} $|f(z)| \leq e^{|z|^k}$ \note{l'exposant est serré et sa lecture douteuse} \ill{}. Alors $\rho$ est la borne inférieure des \struck{\ill{}} nombres $\sigma$ tels que

\page{108}
l'on ait $M(r) \leq e^{r^\sigma}$ pour $r$ assez grand.

\emph{Théorème 2} (Hadamard-Borel) Si l'ordre $\rho$ de $f(z)$ est fini, alors pour \struck{tout} $\varepsilon > 0$, la série
\[
  \text{(5)}\qquad {\sum_n}' \frac{1}{|a_n|^{\rho+\varepsilon}} < +\infty
\]
(l'indice $'$ indique qu'on ne somme pas sur les racines nulles). De plus, on a alors (en supposant $f(0) \neq 0$)
\[
  f(z) = e^{P(z)} \prod_n \Bigl(1 - \frac{z}{a_n}\Bigr)\, e^{\frac{z}{a_n} + \cdots + \frac{z^p}{p a_n^p}}
\]
où $p \leq \rho$, et $P(z)$ est un polynôme de degré $\leq \rho$. \struck{en particulier} Le genre de $f$ est donc $\leq \rho$ \struck{\ill{}} (\uncertain{autrement} \uncertain{dit} genre de $f \leq$ ordre de $f$).

\emph{Corollaire} Si $\rho < 1$, i.e. s'il existe un $\alpha < 1$ tel que $M(r) \leq e^{r^\alpha}$ pour $r$ assez grand, alors $f(z)$ est de genre zéro : $f(z) = C^{\mathrm{te}} \prod (1 - z/a_n)$ \struck{\ill{}}.

\emph{Remarque.} Il est très facile de se convaincre que ordre de $f \leq$ genre de $f + 1$. Donc, si on note $g$ \uncertain{son} genre et $\omega$ l'ordre, on a \struck{\ill{}}
\[
  \omega - 1 \leq g \leq \omega
\]
\note{formule encadrée ; l'ordre, noté $\rho$ jusqu'ici, est désormais $\omega$} Comme $g$ est entier, il en résulte que $g$ est connu dès que $\omega$ l'est, sauf si $\omega$ est entier, \uncertain{auquel} cas on peut avoir $g = \omega$ ou $g = \omega - 1$.

\page{109}
P. ex. $e^{z^\omega}$ \note{l'exposant est surchargé} est une fonction d'ordre $\omega$, de \uncertain{genre} $\omega$. $\prod_n (1 - z/n^\alpha)$ (où $\alpha > 1$) est d'ordre $1/\alpha$ \note{lecture douteuse} (sauf erreur) et de genre $0$.

Notons que l'ordre du produit de deux fonctions entières est au plus égal au \uncertain{maximum} des deux ordres : $\omega(fg) \leq \operatorname{Max}(\omega f, \omega g)$ ; \uncertain{de} \uncertain{même} pour la somme. D'où, \struck{si $f(z) = e^{P(z)}$ \ill{}} \note{en interligne : « l'ensemble des »} fonctions entières dont l'ordre est $< \rho$ (resp. $\leq \rho$) est une algèbre.

\emph{Théorème 3} (Hadamard) Soit la fonction entière $f(z) = \sum c_n z^n$. Si $n^\alpha \sqrt[n]{|c_n|} \to 0$ pour \ill{} \uncertain{fini}, alors $f$ est au plus de genre \ill{} \note{la fin de la ligne sort du champ du cliché ; le feuillet s'arrête là}

\page{110}
\emph{Proposition} Soit $f(z)$ une fonction entière \struck{\ill{}} \note{en interligne : « dont »} les coefficients sont majorés en module par ceux d'une fonction de genre $0$, $h(z)$. Si de plus les racines \uncertain{non} \uncertain{nulles} $z_i$ de $f(z)$ sont telles que $\sum 1/|z_i| < +\infty$, $f(z)$ est de genre zéro.

\struck{En effet, $h(z)$ étant \ill{}} \emph{Démonstration.} On peut supposer que $h(z)$ est de la forme $h(z) = K \prod_{i=1}^{\infty} (1 + \alpha_i z)$, où $\alpha_i \geq 0$, et $\sum \alpha_i < +\infty$. Comme $h(z)$ est d'ordre $\leq 1$, il en est de même de $f(z)$, donc on aura $f(z) = e^{\alpha + \beta z}\, z^m \prod_i (1 - z/z_i)$. Il faut montrer que $\beta = 0$. On peut supposer, \uncertain{en} \uncertain{modifiant} \ill{}, $\beta > 0$, sinon on remplace $f(z)$ par $f(\theta z)$, avec $\theta = \bar{\beta}/|\beta|$ \note{la barre sur $\beta$ est douteuse}. On a alors $|f(z)| = K e^{\beta x}\, |z|^m \prod_i |1 - z/z_i|$ \note{le $K$ porte un exposant « te »}. D'autre part, on sait que \struck{\ill{}} \note{en interligne : « toute »} fonction de genre zéro, telle $h(z)$, admet une majoration de la forme $|h(z)| \leq C^{\mathrm{te}} \times e^{\varepsilon|z|}$, où on peut prendre $\varepsilon$ arbitraire. Si alors on \struck{prend $\varepsilon < \beta$, on aura $|f(z)| \leq |z|^m \prod |1 - z/z_i| \leq K e^{\varepsilon|z|}$} \note{deux lignes biffées}

\page{111}
Notons que si $-\theta \leq \arg z \leq +\theta$, où $\theta$ est un angle arbitraire \ill{} compris \uncertain{strictement} entre \struck{\ill{}} $\pi/3$ et $\pi$\ill{} \note{le second angle est serré en bout de ligne ; $\pi/2$ est ce que la suite demande}, on aura $x \geq c|z|$, où $c$ est une constante $> 0$, d'où $|f(z)| \geq e^{c|z|}\, |z|^m \prod_i |1 - z/z_i|$ quand $-\theta \leq \arg z \leq +\theta$. D'autre part on a $|f(z)| \leq |h(z)| \leq C^{\mathrm{te}} \times e^{\varepsilon|z|}$ d'où, en \uncertain{choisissant} $\varepsilon < c$ : $|z|^m \prod_i |1 - z/z_i| \leq K e^{-\alpha|z|}$ ($\alpha > 0$), dès que $-\theta \leq \arg z \leq +\theta$. Posons $g(z) = z^m \prod_i (1 - z/z_i)$, \uncertain{cette} \uncertain{fonction} \uncertain{satisfait} $|g(z)| \leq K e^{-\alpha|z|}$ si $-\theta \leq \arg z \leq +\theta$ \note{inégalité soulignée}. Montrons qu'une telle inégalité est \uncertain{impossible} pour une fonction $g$ de genre zéro qui n'est pas \struck{\ill{}} $0$. --- \uncertain{On} \uncertain{a} aussi, $g$ étant de genre zéro, $|g(z)| \leq K' e^{\frac{\alpha}{3}|z|}$ pour \ill{}. Posons $\omega = e^{2i\pi/3}$, et $G(z) = g(z)\, g(\omega z)\, g(\omega^2 z)$, c'est une fonction entière, et on a visiblement, \ill{} pour \uncertain{tout} $z$, \ill{} :
\[
  |G(z)| \leq \bigl(K e^{-\alpha|z|}\bigr)\bigl(K' e^{\frac{\alpha}{3}|z|}\bigr)\bigl(K' e^{\frac{\alpha}{3}|z|}\bigr) = K'' e^{-\frac{\alpha}{3}|z|}
\]
$G(z)$ serait donc bornée, donc une constante. Mais $G(z) = z^{3m} \prod_i \Bigl(1 - \dfrac{z^3}{z_i^3}\Bigr)$, il
\marginal{faut donc que l'ensemble des $z_i$ soit vide, et $m = 0$, donc $g(z)$ \uncertain{constante}. \uncertain{Ceci} \uncertain{signifie} $f(z) = C^{\mathrm{te}} e^{\beta z}$, \ill{} ce qui est incompatible avec $|f(z)| \leq e^{k} e^{\varepsilon|z|}$ \ill{}} \note{la note court de bas en haut dans la marge gauche et achève la phrase interrompue au bas de la page. Dans la même marge, à hauteur des premières lignes, un petit croquis : un cercle, et trois rayons issus du centre découpant le disque en trois secteurs, deux des rayons pointillés}

\section{[Ordre d'une suite, fonction de comptage]}

\note{Titre de l'éditeur : les pages 112 à 118 n'en portent pas. Elles poursuivent le même sujet, dans la même écriture, et numérotent (1) et (1') deux formules de la page 115.}

\page{112}
Soit $(\alpha_n)$ une suite décroissante \note{en interligne : « et convergeant vers zéro »} de nombres \uncertain{positifs}. Soit $\rho$ la borne inférieure des nombres $s > 0$ tels que $(\alpha_n)$ soit majorée par un multiple de la suite $(1/n^{1/s})$ (et $+\infty$ si un tel $s$ n'existe pas) et soit $\sigma$ la borne inférieure des nombres $s > 0$ tels que $\sum \alpha_n^s < +\infty$ (et $+\infty$ si un tel $s$ n'existe pas). \uncertain{On} \uncertain{a} aussi $1/\rho = \limsup_n \dfrac{\log(1/\alpha_n)}{\log n}$. \struck{\ill{}} \note{deux lignes biffées} Le critère de Riemann affirme \struck{$\rho \leq \sigma$} $\sigma \leq \rho$. Mais on a aussi $\rho \leq \sigma$, \ill{}

\emph{Lemme} Si $s > 0$ est tel que $\sum_n \alpha_n^s < +\infty$, alors la suite $(\alpha_n)$ est \struck{majorée par un multiple de la suite $(1/n^{1/s})$} $\alpha_n = o(1/n^{1/s})$. ($\alpha_n$ étant décroissant !) \emph{Démonstration.} On se ramène au cas où $s = 1$ (en posant $\beta_n = \alpha_n^s$) et à prouver qu'il est impossible que \struck{\ill{}} pour une suite strictement croissante d'entiers $n_i$, on ait $\alpha_{n_i} \geq 1/n_i$ ; en effet, alors on aurait $\alpha_{n_{i-1}+1} + \alpha_{n_{i-1}+2} + \cdots + \alpha_{n_i} \geq \dfrac{n_i - n_{i-1}}{n_i}$, soit $\varepsilon_i$ ce second membre, $\sum \varepsilon_i$ est \uncertain{donc} convergente, et on a $n_i(1 - \varepsilon_i) = n_{i-1}$, d'où
\[
  n_i = \frac{n_{i-1}}{1 - \varepsilon_i} = \frac{n_{i-2}}{(1 - \varepsilon_i)(1 - \varepsilon_{i-1})} = \cdots = \frac{1}{(1 - \varepsilon_1)\cdots(1 - \varepsilon_i)} \leq \frac{1}{\prod_1^{\infty} (1 - \varepsilon_k)},
\]
\note{le numérateur des deux dernières fractions est lu $1$ ; $n_0$ conviendrait aussi} la suite des $n_i$ est donc majorée, et s'arrête donc \ill{}.

\page{113}
\ill{} Supposons que l'on ait $\sum \alpha_i < +\infty$, $f(z) = \prod_1^{\infty} (1 + \alpha_i z)$. Alors $\rho$ est aussi \ill{} l'ordre de la fonction entière, i.e. $\rho = \limsup_{r \to \infty} \dfrac{\log\log f(r)}{\log r}$. \ill{} En effet, on sait que $\rho$ est $\leq$ ordre de $f$ \ill{}. \ill{} on est \uncertain{ramené} à montrer que : \uncertain{dans} le cas $\alpha_n = 1/n^\beta$, on a $\limsup \dfrac{\log\log f(r)}{\log r} \leq 1/\beta$ (on sait \ill{} que c'est $= 1/\beta$). En effet, on a $f(r) = \prod_1^{\nu} (1 + r/n^\beta) \prod_{\nu+1}^{\infty} (1 + r/n^\beta)$, le deuxième facteur est $\leq e^{R(\nu)\, r}$ où $R(\nu) = \sum_{\nu+1}^{\infty} 1/n^\beta \leq \int_\nu^{\infty} dx/x^\beta = \dfrac{1}{\beta - 1}\cdot\dfrac{1}{\nu^{\beta-1}}$. Prenons $\nu$ tel que $r \cdot \dfrac{1}{\nu^{\beta-1}} \leq r^{1/\beta}$, i.e. $\nu \geq r^{1/\beta}$, par exemple $\nu = [r^{1/\beta}] + 1$. Alors $e^{R(\nu)\, r}$ est une fonction de $r$ d'ordre $\leq 1/\beta$. \uncertain{Il} \uncertain{reste} d'autre part que le produit $\prod_1^{\nu} (1 + r/n^\beta) \leq (1 + r)^\nu = e^{\nu \log(1+r)}$, qui est \uncertain{visiblement} d'ordre $\leq 1/\beta$. Rappelons que si on pose $f(r) = \sum a_n r^n$, \ill{} et $B(r) = \operatorname{Sup}_n (a_n r^n)$, alors on a aussi ordre de $f = \limsup_{r \to \infty} \dfrac{\log\log B(r)}{\log r}$.

\page{114}
Soit $f(z) = \prod_1^{\infty} (1 + \lambda_i z)$ une fonction entière de genre $1$ \note{« genre » souligné}, posons $\alpha_i = |\lambda_i|$, $F(z) = \prod_1^{\infty} (1 + \alpha_i z)$. Alors $f$ et $F$ ont même ordre (qui s'exprime donc aussi en fonction de la suite des $|\lambda_i|$). En effet, l'ordre de $f \leq$ ordre de $F$ (évident). D'autre part, ordre de $F =$ ordre de la suite des $|\alpha_i|$ (voir ci-dessus) $\leq$ ordre de $f$ (th\add{éorème} d'Hadamard).

$\Longrightarrow$ On peut parler de l'ordre d'un ensemble \note{en interligne : « dénombrable »} de nombres complexes, \struck{indépendamment de \ill{}} \uncertain{rangés} en suite, \uncertain{comme} la borne inférieure des $s > 0$ tels que $\sum |z_i|^s < +\infty$ \note{encadré}. Suite \uncertain{à} \uncertain{décroissance} \uncertain{rapide} $=$ \uncertain{suite} d'ordre zéro.

Ordres d'une suite dans un espace de Banach, \ill{} (deux définitions, suivant qu'on \ill{} sommabilité \ill{}, ou sommabilité absolue). Ordre d'une partie bornée (\struck{plus petit} \note{en interligne : « borne inférieure »} des $s > 0$ tels que la partie soit contenue dans l'enveloppe \uncertain{convexe} \uncertain{fermée} d'une suite d'ordre $s$). \note{le bas de la page, une dizaine de lignes au crayon très pâle, n'a pu être lu ; on y distingue des fractions et « $M \geq$ »}

\page{115}
Soit $(r_\nu)$ suite croissante de nombres $> 0$, $\to \infty$, soit $n(r)$ le nombre des indices $\nu$ tels que $r_\nu \leq r$. Soit $0 \leq a \leq b < +\infty$, et $\alpha > 0$, \uncertain{on} \uncertain{a} :
\[
  \text{(1)}\qquad \sum_{\nu,\ a < r_\nu < b} \frac{1}{r_\nu^{\alpha}} = \Bigl(\frac{n(b-0)}{b^{\alpha}} - \frac{n(a+0)}{a^{\alpha}}\Bigr) + \alpha \int_a^b \frac{n(r)\,dr}{r^{\alpha+1}}
\]
(\uncertain{car} le premier membre \uncertain{est} \uncertain{l'intégrale} $\int_a^b r^{-\alpha}\,dn(r)$ \ill{}).

\emph{Proposition} $\sum_\nu 1/r_\nu^{\alpha} < +\infty \iff \int_0^{\infty} \dfrac{n(r)\,dr}{r^{\alpha+1}} < +\infty$, alors $n(r)/r^{\alpha} \to 0$ pour $r \to \infty$, et \uncertain{on} \uncertain{a} :
\[
  \text{(1')}\qquad \sum_{\nu,\ r_\nu > a} \frac{1}{r_\nu^{\alpha}} = -\frac{n(a+0)}{a^{\alpha}} + \alpha \int_a^{\infty} \frac{n(r)\,dr}{r^{\alpha+1}}
\]
\note{formule encadrée} \struck{\ill{}} \note{une ligne biffée} En effet, (1) s'écrit, pour $a = 0$ : $\sum_{r_\nu < b} \dfrac{1}{r_\nu^{\alpha}} = \dfrac{n(b-0)}{b^{\alpha}} + \alpha \int_0^b \dfrac{n(r)\,dr}{r^{\alpha+1}}$ \note{numérotée « (1 bis) », cerclé}, d'où suit que \ill{}. Réciproquement, si $\int_0^{\infty} \dfrac{n(r)\,dr}{r^{\alpha+1}} < +\infty$, alors il existe $M > 0$ et une suite croissante $b_i \to \infty$ tels que $n(b_i)/b_i^{\alpha} \leq M$ ; \ill{} on \uncertain{aurait} $\sum_{r_\nu < b_i} \dfrac{1}{r_\nu^{\alpha}} \leq M + \alpha \int_0^{\infty} \dfrac{n(r)\,dr}{r^{\alpha+1}}$, d'où $\sum_\nu 1/r_\nu^{\alpha} < +\infty$ \note{encadré}.

\page{116}
Montrons que par ailleurs $n(r)/r^{\alpha} \to 0$ ; \uncertain{sinon} $n(b_i)/b_i^{\alpha} > \varepsilon$ pour une suite $b_i \to +\infty$, \uncertain{et} on aurait $n(b_i) - n(b_{i-1}) \geq \varepsilon \lambda_i^{\alpha} - M \lambda_{i-1}^{\alpha}$ \note{ainsi sur la page : la suite est notée $b_i$ dans le texte et $\lambda_i$ dans les formules}, d'où $\sum_{b_{i-1} < r_\nu \leq b_i} r_\nu^{-\alpha} \geq (n(b_i) - n(b_{i-1}))\,\lambda_i^{-\alpha} \geq \varepsilon - M(\lambda_{i-1}/\lambda_i)^{\alpha}$, \struck{\ill{}} \uncertain{et} \uncertain{on} \uncertain{peut} \uncertain{supposer} $M(\lambda_{i-1}/\lambda_i)^{\alpha} \leq \varepsilon/2$ pour tout $i$, d'où $\sum_{b_{i-1} < r_\nu \leq b_i} r_\nu^{-\alpha} \geq \varepsilon/2$, ce qui est contradictoire. Si alors on fait $b \to \infty$ dans (1), on obtient (1'). \struck{\ill{}} Des résultats analogues \uncertain{s'obtiennent} \uncertain{si} on considère \ill{} les sommes $\sum_\nu \lambda(r_\nu)$, où $\lambda$ est une fonction définie dans $\mathbf{R}^{+*}$, continûment dérivable, \uncertain{décroissante} et tendant vers zéro. Alors :
\[
  \sum_\nu \lambda(r_\nu) < +\infty \iff \int_0^{\infty} n(r)\,(-\lambda'(r))\,dr < +\infty
\]
\note{encadré} \struck{\ill{}} et alors $\lambda(r)\,n(r) \to 0$ pour $r \to \infty$, et
\[
  \sum_{r_\nu > a} \lambda(r_\nu) = -\lambda(a)\,n(a+0) + \int_a^{\infty} n(r)\,(-\lambda'(r))\,dr
\]
\note{encadré}

\page{117}
Supposons maintenant que $(r_\nu)$ soit la suite des modules des zéros d'une fonction entière $f(z)$. On a $2\pi\, n(r) = \int_0^{2\pi} \dfrac{f'(re^{i\varphi})}{f(re^{i\varphi})}\, re^{i\varphi}\, d\varphi$, d'où
\[
  2\pi \int_a^b n(r)\,\mu(r)\,dr = \int_a^b \int_0^{2\pi} \frac{f'(re^{i\varphi})}{f(re^{i\varphi})}\, e^{i\varphi}\, r\mu(r)\, dr\, d\varphi
\]
\[
  = \int_0^{2\pi} d\varphi \int_a^b \frac{f'(re^{i\varphi})}{f(re^{i\varphi})}\, e^{i\varphi}\, r\mu(r)\, dr = \int_0^{2\pi} d\varphi \int_a^b \frac{d}{dr}\log f(re^{i\varphi})\; r\mu(r)\, dr
\]
\[
  = \int_0^{2\pi} d\varphi \Bigl( \bigl[ r\mu(r) \log f(re^{i\varphi}) \bigr]_a^b - \int_a^b \log f(re^{i\varphi})\, d(r\mu(r)) \Bigr)
\]
(en supposant $\mu(r)$ continûment différentiable). --- On devrait bien entendu \uncertain{préciser} la détermination du logarithme \ill{} ; \struck{simplement} d'ailleurs, \ill{} on \ill{}, en prenant les parties réelles, \ill{} : \note{le passage aux parties réelles tient en deux lignes en grande partie illisibles}
\[
  \int_a^b n(r)\,\mu(r)\,dr = \bigl[ r\mu(r)\,V(r) \bigr]_a^b - \int_a^b V(r)\, d(r\mu(r))
\]
\note{encadré ; un « $2\pi$ » biffé en tête du premier membre} où $V(r) = \dfrac{1}{2\pi}\int_0^{2\pi} \log|f(re^{i\varphi})|\,d\varphi$. Prenons $\mu(r) = -\lambda'(r)$, on obtient ($0 < a < b < +\infty$) :
\[
  \sum_{a < r_\nu < b} \lambda(r_\nu) = \lambda(b)\,n(b-0) - \lambda(a)\,n(a+0) + \bigl[ r\mu(r)\,V(r) \bigr]_a^b + \int_a^b V(r)\, d(r\lambda'(r))
\]
\note{« $+\int$ » biffé avant le dernier terme, puis récrit}
\marginal{\uncertain{a} \uncertain{priori} \ill{}} \note{deux mots obliques dans la marge gauche, en regard du milieu de la page, et un cercle tracé à main levée au-dessous}

\page{118}
(Si $r^2\mu(r) \to 0$ pour $r \to 0$, et $|f(0)| = 1$, on a : $\sum_{r_\nu < b} \lambda(r_\nu) = \lambda(b)\,n(b-0) - b\lambda'(b)\,V(b) + \int_0^b V(r)\, d(r\lambda'(r))$.) Comme $|V(r)| \leq M(r) = \operatorname{Sup}_{|z|=r} |f(z)|$ \note{ainsi sur la page ; $\log M(r)$ est ce que $V$ demande}, on obtient une majoration de \struck{$\sum \lambda(r_\nu)$} $\int_a^b n(r)\,\mu(r)\,dr$ en fonction de $M(r)$. \note{le reste de la page est blanc}

\section{Sur une certaine fonction holomorphe sur l'espace de Banach $\ell^1$}

\note{Le titre est de lui, souligné, en tête de la page 119 ; « définie » y est biffé entre « holomorphe » et « sur ».}

\page{119}
Soit, pour $\lambda = (\lambda_i) \in \ell^1$ et tout entier $n \geq 0$ :
\[
  \alpha_n(\lambda) = \sum_{i_1 < \cdots < i_n} \lambda_{i_1} \cdots \lambda_{i_n}
\]
\[
  \alpha_n(\lambda) \leq \frac{\|\lambda\|_1^{\,n}}{n!}
\]
Les $\alpha_n(\lambda)$ sont les coefficients de la fonction entière de genre zéro $f(z) = \prod (1 + \lambda_i z)$, et on obtient \ill{} \uncertain{les} coefficients de \uncertain{toutes} les fonctions \uncertain{entières} de genre zéro.

\emph{Proposition} Pour tout $\lambda \in \ell^1$, la série $\sum_0^{\infty} n!\,\alpha_n(\lambda)$ est absolument convergente, soit $S(\lambda)$ sa somme. La fonction $\lambda \mapsto S(\lambda)$ est holomorphe au voisinage de $0$, avec un rayon de convergence égal à $1$, et ses restrictions \ill{} sont holomorphes. \note{la fin de l'énoncé est en partie perdue ; une accolade verticale réunit ses lignes, une autre celles du corollaire}

\emph{Corollaire} Pour tout $\lambda \in \ell^1$, \ill{}
\[
  \sqrt[n]{|\alpha_n(\lambda)|} = o\Bigl(\frac{1}{n}\Bigr),
\]
\uncertain{donc} si $\sum a_n z^n$ est une fonction entière de genre zéro, \ill{} $\sqrt[n]{n!\,|\alpha_n(\lambda)|} \to 0$, \uncertain{i.e.} $\alpha_n(\lambda) = \dfrac{\varepsilon_n^{\,n}}{n!}$ avec $\varepsilon_n \to 0$, \ill{} \uncertain{une} \uncertain{fonction} \uncertain{entière} \uncertain{et} \ill{}. En effet, $\sqrt[n]{n!\,|\alpha_n(\lambda)|} \to 0$ \uncertain{exprime} que

\page{120}
$\sum n!\,\alpha_n(\lambda z) = \sum n!\,\alpha_n(\lambda)\, z^n$ \uncertain{converge} quel que soit $z$. Les deux relations \struck{\ill{}} \uncertain{sont} \uncertain{les} \uncertain{mêmes} \ill{}, \uncertain{la} \uncertain{seconde} \uncertain{exprimant} \uncertain{la} \uncertain{première} \ill{} (Stirling : $\sqrt[n]{n!} \sim \dfrac{n}{e}$).

\emph{Démonstration} \uncertain{de} \uncertain{la} \uncertain{proposition} \note{« Démonstration » est souligné ; la ligne est biffée d'un trait puis reprise}. \uncertain{On} \uncertain{peut} \uncertain{supposer} \uncertain{les} $\lambda_i > 0$ \uncertain{et} \uncertain{décroissants}. Soit alors $f(z) = \prod (1 + \lambda_i z) = \sum \alpha_n z^n$, où $\alpha_n = \alpha_n(\lambda)$ \ill{}. \uncertain{Prenons} $R > 0$ \ill{}, \uncertain{alors} \uncertain{Cauchy} : $\alpha_n \leq \dfrac{M(R)}{R^{n+1}}$ \note{ainsi sur la page ; $R^n$ attendu}, où $M(R)$ est \uncertain{le} \uncertain{maximum} de $|f(z)|$ sur la circonférence $|z| = R$, donc $M(R) = \prod (1 + \lambda_i R)$. \struck{Soit} \uncertain{Or} \ill{} $\prod_{i=n+2}^{\infty} (1 + \lambda_i R) \leq e^{\rho_{n+1} R}$, \uncertain{où} $\rho_{n+1} = \sum_{k=n+2}^{\infty} \lambda_k$. \uncertain{Prenons} $R = \dfrac{n+1}{\rho_{n+1}}$, \uncertain{on} \uncertain{obtient} :
\[
  \alpha_n \leq \Bigl(\frac{\rho_{n+1}}{n+1} + \lambda_1\Bigr)\cdots\Bigl(\frac{\rho_{n+1}}{n+1} + \lambda_{n+1}\Bigr) e^{n+1} \leq \Bigl(\frac{\rho_1}{1} + \lambda_1\Bigr)\Bigl(\frac{\rho_2}{2} + \lambda_2\Bigr)\cdots\Bigl(\frac{\rho_{n+1}}{n+1} + \lambda_{n+1}\Bigr) e^{n+1}
\]
\note{les exposants de $e$ sont lus $n$ ou $n+1$ ; un facteur intermédiaire est biffé} Or, \struck{\ill{}} \uncertain{les} $\lambda_i$ \uncertain{formant} \uncertain{une} \uncertain{suite} \uncertain{décroissante}, \uncertain{et} \uncertain{comme} $\lambda_n = o(1/n)$, \uncertain{d'autre} \uncertain{part} \uncertain{aussi} $\dfrac{\rho_n}{n} = o\Bigl(\dfrac{1}{n}\Bigr)$, d'où $\dfrac{\rho_n}{n} + \lambda_n \leq \dfrac{\varepsilon_n}{n}$, \uncertain{où} $\varepsilon_n \to 0$. D'où
\[
  \alpha_n \leq \frac{\varepsilon_0\, \varepsilon_1 \cdots \varepsilon_{n+1}}{(n+1)!}\, e^{n+1}
\]
\ill{} \uncertain{où} $\varepsilon_i \to 0$. \uncertain{D'où} \ill{} \note{la page s'arrête là ; la suite est au lot suivant}

\end{document}
